\chapter{TINJAUAN PUSTAKA}

\section{Penyimpangan Konfigurasi (\textit{Configuration Drift})}

Penyimpangan konfigurasi (\textit{configuration drift}) adalah ketidakcocokan antara keadaan aktual sistem dan keadaan yang ditetapkan sebagai acuan deklaratif \citep{pohjola2025mitigating}. Drift berbeda dari cacat statis: konfigurasi yang tidak aman dapat tetap cocok dengan acuannya, sedangkan perubahan aktual yang tidak cocok dengannya merupakan drift meskipun nilai aktual belum dinilai dampaknya. Penelitian ini mengukur kesetaraan terhadap deklarasi, bukan kebenaran semantik atau keamanan seluruh manifest.

Penelitian empiris oleh \citet{zhang2026configuration} melakukan studi terhadap 719 \textit{defect} konfigurasi dari 2.260 skrip Kubernetes dan mengidentifikasi 15 kategori \textit{defect}, di mana 533 dari 719 \textit{defect} (74,1\%) menyebabkan \textit{crash}, operasi yang salah, atau \textit{outage}. Studi tersebut menegaskan bahwa \textit{misconfiguration} merupakan risiko operasional utama pada Kubernetes. \citet{rahman2023security} menambahkan bahwa \textit{security misconfiguration} pada manifest Kubernetes \textit{open-source} mencakup 11 kategori yang berulang, termasuk \textit{container} yang berjalan sebagai \textit{root}, tanpa \textit{drop capabilities}, dan tanpa \textit{resource limits}. Penelitian menunjukkan bahwa penyimpangan konfigurasi bukan hanya masalah estetika konfigurasi, tetapi ancaman nyata terhadap keandalan dan keamanan layanan.

\citet{pohjola2025mitigating} memberikan kerangka konseptual untuk memahami sumber dan dampak penyimpangan pada sistem \textit{Infrastructure-as-Code} (IaC). Penelitian tersebut mengidentifikasi tiga faktor pendorong utama penyimpangan: (1)~keahlian manusia (\textit{human expertise}) yang tidak konsisten; (2)~tekanan waktu (\textit{time pressure}) yang memicu perubahan cepat dan tidak didokumentasikan; dan (3)~ketiadaan mekanisme deteksi otomatis pada \textit{pipeline deployment}. Meskipun penelitian tersebut tidak beroperasi pada Kubernetes secara eksklusif, prinsip-prinsip yang dikemukakan berlaku umum untuk platform orkestrasi berbasis konfigurasi deklaratif.

\subsection{Faktor yang Menimbulkan Penyimpangan Konfigurasi}

Berdasarkan literatur, faktor-faktor utama yang memperburuk penyimpangan konfigurasi pada Kubernetes meliputi:

\begin{enumerate}
  \item Perubahan langsung pada klaster: Penggunaan \texttt{kubectl\allowbreak{} edit}, \texttt{kubectl\allowbreak{} patch}, atau \texttt{kubectl\allowbreak{} scale} secara langsung tanpa memperbarui repositori \textit{source of truth} \citep{zhang2026configuration}.
  
  \item Prosedur \textit{deployment} yang tidak terdokumentasi: Setiap \textit{operator} dapat memiliki urutan perintah yang berbeda, yang mengakibatkan \textit{state} klaster tidak deterministik \citep{pohjola2025mitigating}.
  
  \item Kurangnya mekanisme \textit{rollback} yang terintegrasi: kemampuan untuk mengembalikan perubahan bergantung pada pencatatan versi, ketersediaan nilai dasar, dan fitur yang dikonfigurasi pada alat pengelolaan \citep{rahman2023security}.
  
  \item Ketiadaan mekanisme deteksi penyimpangan otomatis: Tanpa pemeriksaan berkala antara \textit{state} aktual dan \textit{desired state}, penyimpangan dapat terakumulasi selama periode lama tanpa terdeteksi \citep{zhang2026configuration}.
\end{enumerate}

\section{Terminologi Dasar}

Bagian ini mendefinisikan terminologi kunci yang digunakan dalam penelitian, dengan tujuan membangun dasar konseptual yang konsisten sebelum pembahasan lebih lanjut.

\subsection{Kubernetes dan Orkestrasi Kontainer}
Kubernetes adalah platform \textit{container orchestration} \textit{open-source} yang mengotomasi \textit{deployment}, \textit{scaling}, dan manajemen aplikasi terkemas (\textit{containerized}) \citep{burns2016borg}. Kubernetes awalnya dikembangkan oleh Google berdasarkan pengalaman mengelola sistem \textit{workload} berskala besar di Borg dan Omega, dan kini menjadi standar de facto untuk orkestrasi kontainer di industri \citep{verma2015borg}. Pada Kubernetes, \textit{state} yang diinginkan (\textit{desired state}) dideklarasikan melalui berkas konfigurasi (\textit{manifest}), dan \textit{control plane} bertanggung jawab mencapai dan mempertahankan \textit{state} tersebut.

\subsection{Git dan \textit{Version Control}}
Git adalah sistem \textit{version control} terdistribusi yang merekam setiap perubahan pada berkas beserta metadata (penulis, \textit{timestamp}, pesan \textit{commit}) \citep{chacon2014git}. Dalam konteks GitOps, Git berfungsi sebagai \textit{single source of truth} yang menyimpan deklarasi \textit{state} infrastruktur dan aplikasi, sehingga setiap perubahan dapat dilacak, di-\textit{review}, dan dikembalikan (\textit{rollback}) secara andal.

\subsection{\textit{Continuous Integration} dan \textit{Continuous Deployment}}
\textit{Continuous Integration} (CI) adalah praktik di mana perubahan kode secara otomatis diuji dan digabungkan ke \textit{branch} utama, sementara \textit{Continuous Deployment} (CD) adalah perluasan CI di mana perubahan yang lulus pengujian secara otomatis di-\textit{deploy} ke lingkungan produksi tanpa intervensi manual eksplisit \citep{humble2010continuous}. GitOps merupakan paradigma spesifik di dalam CD yang menggunakan repositori Git sebagai \textit{single source of truth}.

\subsection{\textit{Infrastructure-as-Code} (IaC)}
\textit{Infrastructure-as-Code} (IaC) adalah praktik mendefinisikan dan mengelola infrastruktur melalui berkas konfigurasi yang diberi versi (\textit{versioned}), bukan melalui proses manual \citep{morris2020iac}. IaC memungkinkan \textit{version control}, \textit{code review}, dan reprodusibilitas lingkungan. Dalam konteks Kubernetes, IaC diimplementasikan melalui \textit{manifest} deklaratif (YAML) yang mendefinisikan \textit{desired state} klaster.

\subsection{Paradigma Deklaratif vs Imperatif}
Dalam paradigma imperatif, tindakan pengelolaan dinyatakan sebagai perintah atau urutan tindakan yang mengubah keadaan sistem (\textit{how to achieve the state}). Dalam paradigma deklaratif, keadaan yang diinginkan dinyatakan dan pengendali bertugas mendekatinya (\textit{what state to achieve}) \citep{scott2015plp}. Keduanya dapat diotomasi; perbedaan untuk penelitian ini adalah prosedur koreksi imperatif tetap yang dijalankan melalui \texttt{kubectl} dibandingkan dengan rekonsiliasi ArgoCD.

\section{Manajemen Konfigurasi Infrastruktur}

\textit{Infrastructure state management} mengatur bagaimana keadaan infrastruktur didefinisikan, diterapkan, dan dipertahankan. Pendekatan deklaratif menyimpan keadaan yang diinginkan dalam konfigurasi, sedangkan pendekatan imperatif menyatakan tindakan untuk mengubah keadaan \citep{morris2020iac, limoncelli2018}. Otomasi, pemeriksaan berulang, dan pencatatan dapat tersedia pada kedua pendekatan; karakteristik tersebut perlu ditetapkan pada tingkat implementasi yang dibandingkan.

\section{\textit{Continuous Deployment} dan Paradigma GitOps}

\textit{Continuous Deployment} (CD) adalah praktik dalam ranah DevOps di mana setiap perubahan kode yang lulus pengujian secara otomatis di-\textit{deploy} ke lingkungan produksi tanpa intervensi manual eksplisit \citep{limoncelli2018}. CD merupakan perluasan dari \textit{Continuous Integration} (CI) dan bertujuan meminimalkan waktu antara penulisan kode dan ketersediaannya bagi pengguna akhir. Dalam ekosistem Kubernetes, CD memerlukan mekanisme yang dapat secara andal menyelaraskan perubahan konfigurasi dengan \textit{state} klaster.

GitOps adalah paradigma spesifik di dalam CD yang menggunakan repositori Git sebagai \textit{single source of truth} untuk mendefinisikan \textit{state} infrastruktur dan aplikasi. CNCF GitOps Working Group mendefinisikan empat prinsip OpenGitOps: (1)~sistem dideklarasikan secara deklaratif; (2)~\textit{state} sistem yang diinginkan tersimpan dan diberi versi (\textit{versioned}) dalam Git; (3)~perubahan \textit{state} ditarik (\textit{pulled}) secara otomatis oleh agen \textit{software}; dan (4)~agen \textit{software} secara kontinu mendeteksi dan mengoreksi deviasi dari \textit{state} yang diinginkan \citep{opengitops2021}.

GitOps mendukung pencatatan versi deklarasi melalui Git, peninjauan perubahan, dan penggunaan revisi sebelumnya sebagai acuan pemulihan \citep{yuen2021gitops}. Riwayat Git tidak otomatis merekam perubahan langsung pada klaster; tindakan tersebut memerlukan bukti dari pengendali atau audit API. Pada model \textit{pull-based}, agen di dalam klaster dapat menarik deklarasi tanpa pipeline eksternal menerapkan setiap perubahan ke API klaster, tetapi kebutuhan kredensial dan akses operasional tetap bergantung pada rancangan serta konfigurasi sistem \citep{utami2021analysis}. Klaim adopsi dan manfaat GitOps juga muncul dalam artikel industri CNCF 2025; artikel tersebut dipakai hanya sebagai konteks industri, bukan sebagai bukti \textit{peer-review} \citep{cncf2025gitops}.

\subsection{Performa \textit{Deployment} dan \textit{Anti-Pattern}}

Forsgren dkk. \citep{forsgren2018accelerate} membahas \textit{deployment frequency}, \textit{lead time for changes}, \textit{time to restore service}, dan \textit{change failure rate} sebagai metrik performa pengembangan dan operasi. Metrik organisasi tersebut tidak identik dengan pengukuran eksperimen ini dan tidak dengan sendirinya membuktikan bahwa satu prosedur selalu lebih unggul. Penelitian ini menggunakan konstruk yang lebih terbatas: kesetaraan bidang konfigurasi, waktu deteksi dan pemulihan konfigurasi, serta bukti keterlacakan.

ArgoCD menyediakan pengendali aplikasi yang membandingkan sumber daya teramati dengan deklarasi repositori dan, jika kebijakan serta cakupannya mengizinkan, dapat melakukan sinkronisasi atau koreksi otomatis \citep{argocd2024}. Perilaku aktual bergantung pada konfigurasi aplikasi, izin, ketersediaan pengendali, jenis sumber daya, dan cakupan inventaris. Oleh sebab itu, rekonsiliasi tidak menjamin setiap perubahan akan selalu ditemukan atau dipulihkan.

\section{Dasar Teori: \textit{Reconciliation} dan \textit{Convergence} pada Sistem Deklaratif}

Karakteristik GitOps yang berbeda dari pendekatan imperatif dapat dijelaskan melalui dua prinsip teoretis yang saling berkaitan.

\subsection{\textit{Convergence} Menuju \textit{Desired State}}

Pengendali deklaratif secara berkala membandingkan keadaan aktual dengan keadaan yang diinginkan dan dapat melakukan tindakan korektif sesuai kebijakan yang berlaku \citep{argocd2024}. Keberhasilan konvergensi bukan sifat yang dapat diasumsikan untuk semua sumber daya atau gangguan: hasil dipengaruhi oleh cakupan aplikasi, izin, ketersediaan komponen, konfigurasi sinkronisasi, dan sifat sumber daya.

Perintah imperatif tidak dengan sendirinya menyediakan rekonsiliasi berkelanjutan. Namun, prosedur imperatif dapat diotomasi dengan pemeriksaan berkala dan tindakan koreksi yang telah ditentukan. Karena penelitian ini membandingkan prosedur \texttt{kubectl} terotomasi dengan ArgoCD, bukan respons manusia, perbedaan waktu tidak ditafsirkan sebagai perbedaan kecepatan operator.

\subsection{\textit{Audit Trail} sebagai \textit{Social Control Mechanism}}

Riwayat Git dapat menyediakan metadata dan perubahan deklarasi, sedangkan log pengendali dan audit API dapat memberi bukti tambahan tentang penerapan aktual \citep{opengitops2021}. Riwayat tersebut tidak otomatis membuktikan siapa yang melakukan perubahan pada klaster atau apakah semua tindakan terekam; kelengkapan bergantung pada tata kelola repositori, identitas, konfigurasi audit, dan pengumpulan bukti. Namun, R3 tidak diukur dalam penelitian ini karena audit API server Kubernetes tidak tersedia; keterlacakan dibahas sebagai keterbatasan dan pekerjaan lanjutan, bukan sebagai hasil.

\section{\textit{Runbook} sebagai Standarisasi Prosedur Pemulihan}

\textit{Runbook} adalah dokumen terstandarisasi yang berisi prosedur operasional untuk menanggapi insiden atau melakukan pemulihan (\textit{recovery}) pada sebuah sistem. Dalam ranah operasi TI, \textit{runbook} berfungsi sebagai \textit{playbook} yang mendefinisikan langkah-langkah korektif berurutan berdasarkan jenis insiden, sehingga setiap \textit{operator} yang berbeda mengikuti urutan tindakan yang sama \citep{schlette2024playbook}. Penggunaan \textit{runbook} mengurangi variabilitas akibat keahlian \textit{operator} yang berbeda-beda dan memastikan bahwa keputusan subjektif diminimalkan selama pemulihan.

Dalam rancangan penelitian ini, prosedur imperatif bukan runbook yang dipandu kepada operator. Mode imperatif menggunakan pemeriksa dan korektor otomatis yang tetap, dengan pembandingan deklarasi dan keadaan aktual serta tindakan \texttt{kubectl} yang telah ditetapkan. Mode ini diperoleh dengan menjeda \textit{auto-sync}/\textit{self-heal} ArgoCD pada stack yang sama, lalu memulihkannya. Pilot perlu memverifikasi bahwa prosedur tersebut bekerja tanpa menerima identitas skenario atau label injeksi dan menyimpan bukti yang sama jenisnya dengan evaluator bersama.

\section{Paradigma \textit{Deployment} Alternatif}

Pilihan pengelolaan konfigurasi mencakup rekonsiliasi berbasis tarik, pipeline dorong, alat konfigurasi yang dijalankan sesuai kebutuhan, serta otomasi imperatif khusus. Perbedaan autentikasi, cakupan sumber daya, pencatatan, dan frekuensi eksekusi bergantung pada konfigurasi implementasi; tidak diasumsikan bahwa seluruh pipeline dorong mengekspos kredensial atau bahwa semua alat imperatif tidak memiliki audit. Studi Shrestha dan Ali secara khusus membandingkan GitOps dengan Ansible pada waktu induksi, penemuan, dan pemulihan drift \citep{shrestha2024configuration}.

Penelitian ini membatasi pembanding pada satu prosedur imperatif terotomasi berbasis \texttt{kubectl} dan satu instalasi ArgoCD. Kontras ini tidak merupakan batas atas efek yang berlaku bagi semua alat atau konteks; temuan nantinya hanya berlaku bagi implementasi dan workload yang benar-benar diuji.

\section{Evaluasi Empiris GitOps dan ArgoCD}

Studi terkait telah menguji sejumlah kombinasi alat dan tujuan penelitian. Shrestha dan Ali \citep{shrestha2024configuration} membandingkan GitOps dengan Ansible pada waktu induksi, penemuan, dan pemulihan drift. Kagganti Nataraja \citep{kaggantinataraja2025security} mendeskripsikan perbandingan terkontrol pada klaster Minikube satu node dengan konfigurasi keamanan yang berbeda. Matubber \citep{matubber2025managing} mengevaluasi pengelolaan infrastruktur Kubernetes berbasis GitOps dan mencakup pembanding lama yang semi-terotomasi. Ketiganya menunjukkan bahwa penelitian komparatif atau evaluasi berbasis waktu telah ada; abstrak yang tersedia tidak cukup untuk menyimpulkan ada atau tidaknya uji statistik tertentu.

Utami dan Al Fatta \citep{utami2021analysis} membahas model deployment deklaratif berbasis tarik dibandingkan dengan model berbasis dorong. Paavola \citep{paavola2021managing} membahas pengelolaan beberapa aplikasi Kubernetes menggunakan prinsip GitOps, sedangkan D'Amore \citep{damore2021gitops} mendeskripsikan implementasi ArgoCD pada lingkungan Kubernetes hybrid cloud. Studi tersebut memberi konteks implementasi, tetapi tidak diperlakukan di sini sebagai bukti bahwa salah satu prosedur selalu lebih unggul.

Perbandingan penelitian perlu dibaca sesuai tujuan, alat, lingkungan, konfigurasi, dan hasil yang benar-benar dilaporkan oleh masing-masing sumber. Draf ini tidak menyatakan bahwa penelitian terdahulu tidak terkontrol, tidak menggunakan statistik, atau tidak memiliki baseline apabila pernyataan itu tidak dapat dipastikan dari abstrak yang tersedia.

\section{Metode Eksperimental Terkontrol pada Riset Sistem}

Evaluasi empiris pada riset sistem dan infrastruktur memerlukan desain yang memitigasi ancaman terhadap validitas \citep{wohlin2012experimentation}. Penelitian ini merupakan studi kasus komparatif terkontrol: satu kasus diamati pada kondisi dasar, satu intervensi, dan kondisi pascaintervensi, dengan dua mode pengelolaan yang dijalankan berurutan \citep{runeson2009case}. \citet{limoncelli2018} menjelaskan bahwa penelitian sistem harus membangun lingkungan yang dapat direproduksi, mengendalikan variabel pengganggu, dan mendefinisikan metrik yang jelas sebelum menyimpulkan hubungan kausal.

Rancangan penelitian ini menetapkan mode imperatif terotomasi berbasis \texttt{kubectl} dan mode GitOps berbasis ArgoCD sebagai dua perlakuan pada satu stack. Kesetaraan \textit{workload}, nilai dasar, inventaris pengukuran, dan prosedur injeksi direncanakan untuk mengurangi perbedaan yang tidak relevan. Karena kedua implementasi dapat berbeda dalam lebih dari satu komponen teknis dan tidak ada isolasi penuh, hasil tidak langsung diatribusikan sebagai efek universal paradigma; interpretasinya dibatasi pada prosedur dan kasus yang benar-benar diuji \citep{wohlin2012experimentation}.

Analisis yang direncanakan bersifat per skenario dan eksploratif. Ringkasan distribusi, jumlah observasi lengkap dan tersensor, serta ukuran efek dan interval kepercayaan dilaporkan bersama-sama; uji komparatif hanya digunakan jika asumsi dan struktur data mendukungnya. Sampel formal 15 pengulangan per skenario per mode (210 run) tidak dianggap sebagai bukti daya statistik yang memadai \citep{dyba2006power,kitchenham2024small}. Rincian definisi operasional dan batas interpretasi terdapat pada Bab III.

\section{Klasifikasi Beban Uji (\textit{Workload})}

Aplikasi yang dijalankan pada klaster Kubernetes dapat dikelompokkan ke dalam beberapa kelas \textit{workload} berdasarkan karakteristik arsitektur dan manajemen \textit{state} \citep{siddiqui2023microservices, aghili2024workload}:

\begin{enumerate}
  \item \textit{Workload} web tanpa state (\textit{stateless web microservice}): Layanan yang tidak menyimpan \textit{state} secara internal, sehingga setiap permintaan dapat dilayani oleh instans mana pun. Contoh: server HTTP sederhana, \textit{API gateway}.
  
  \item \textit{Workload} berbasis data yang menyimpan \textit{state} (\textit{stateful data service}): Layanan yang menyimpan \textit{state} secara persisten, seperti \textit{database} relasional, \textit{search engine}, atau \textit{object storage}. \textit{Workload} kelas ini memerlukan urutan \textit{deployment} dan penanganan penyimpangan yang lebih kompleks.
  
  \item \textit{Workload} komposit multi-service: Aplikasi yang terdiri atas banyak komponen yang saling bergantung, mencakup kombinasi kelas 1 dan 2 beserta komponen infrastruktur (\textit{ingress}, TLS, \textit{secrets}, \textit{monitoring}, \textit{backup}, kebijakan keamanan). Kelas \textit{workload} ini memiliki kompleksitas dependensi tertinggi dan paling rentan terhadap penyimpangan konfigurasi akibat banyaknya komponen yang harus diselaraskan.
\end{enumerate}

Pengujian formal dalam rancangan ini menggunakan beban uji komposit multi-service InvenioRDM pada satu stack produksi-like. Inventaris dan komponen aktual harus dicatat dari konfigurasi yang divalidasi sebelum eksperimen; jumlah komponen yang pernah disebut sebagai target belum diverifikasi. Beban uji merupakan konteks pengukuran, bukan dasar untuk menggeneralisasi hasil ke seluruh \textit{workload} multi-service. Surrogat S1/S2 historis/offline tidak dipakai dalam desain formal dan tidak memperluas sampel formal.

\section{Ruang Kontribusi Penelitian}

Studi komparatif mengenai GitOps dan pengelolaan konfigurasi telah ada, termasuk evaluasi terhadap Ansible, pengujian terkontrol pada Minikube, dan pengukuran yang mencakup \textit{baseline} semi-terotomasi. Karena itu, penelitian ini tidak berangkat dari klaim bahwa perbandingan terdahulu tidak ada atau bahwa keunggulan GitOps telah terbukti secara universal.

Kontribusi yang diklaim bersifat terbatas dan dapat diperiksa: (1)~perbandingan empiris pada \textit{deployment} InvenioRDM multi-service nyata, bukan \textit{workload} sintetis; (2)~dua mode yang dijalankan berurutan pada satu stack produksi-like, dengan mode imperatif didefinisikan sebagai penjedaan ArgoCD ditambah pemeriksa/\textit{runbook} \texttt{kubectl}, sehingga perbandingan tidak memerlukan stack kembar; (3)~pengukuran R1/R2 yang dapat direproduksi dengan kepemilikan HPA, \textit{drift} Service F yang dilaporkan terpisah, dan kasus G yang dipisahkan dari efikasi; serta (4)~bukti batas bahwa hasilnya berlaku untuk kasus dan apparatus yang diuji, bukan klaim superioritas GitOps universal. R3 tidak diukur karena audit API server tidak tersedia dan dibahas sebagai keterbatasan serta pekerjaan lanjutan.

Hasil yang kelak diperoleh hanya dapat menjawab pertanyaan pada skenario, \textit{workload}, konfigurasi, dan resolusi instrumentasi yang benar-benar diuji. Besaran maupun arah perbedaan belum dapat dipastikan sebelum pengumpulan serta analisis data formal.
